\chapter{LIMITAÇÕES E AMEAÇAS À VALIDADE}
\label{chap:limitacoes}

Este capítulo delimita o alcance das evidências antes de detalhar cada ameaça à validade. A Tabela~\ref{tab:limitacoes-sintese} organiza as limitações centrais por tipo, explicita sua consequência e indica a principal estratégia de mitigação.

\begin{table}[H]
  \centering
  \caption{Síntese das limitações, consequências e estratégias de mitigação}
  \label{tab:limitacoes-sintese}
  \small
  \renewcommand{\arraystretch}{1.12}
  \begin{tabularx}{\textwidth}{@{}>{\raggedright\arraybackslash}p{2.45cm}XXX@{}}
    \toprule
    \textbf{Tipo} & \textbf{Limitação central} & \textbf{Consequência} & \textbf{Mitigação prioritária} \\
    \midrule
    Interna & Um particionamento fixo e critérios de precisão definidos sobre as perdas médias. & Sensibilidade ao conjunto de teste e incerteza não quantificada diretamente nos indicadores de preservação. & Repetir particionamentos válidos e estimar intervalos dos indicadores por reamostragem. \\
    \addlinespace
    Construto & $\rankrho$, $\AW$, $\AR$ e $\SR$ representam componentes distintos e usam perda 0--1. & Nenhum indicador isolado descreve todo o perfil de cooperação preditiva ou a utilidade da decisão. & Interpretar o perfil conjunto e avaliar métricas preditivas e custos alternativos. \\
    \addlinespace
    Externa & Três datasets binários, tabulares e de baixa dimensão; Random Forest apenas no SUPPORT2. & Os resultados não se generalizam automaticamente a outros domínios ou classificadores. & Ampliar datasets, dimensionalidades, classes e famílias de modelos. \\
    \addlinespace
    Computacional & A enumeração cresce como $2^d-1$. & O custo restringe o protocolo exaustivo a poucos atributos. & Avaliar busca estruturada ou amostragem de subconjuntos sem perder comparações relevantes. \\
    \addlinespace
    Extrapolação & A grade termina no maior $n$ avaliado em cada dataset. & A última inversão observada pode não ser a última fora da grade. & Validar previsões em tamanhos amostrais deliberadamente ocultados. \\
    \bottomrule
  \end{tabularx}
  \sourcebelow{elaboração própria}
\end{table}

\section{Validade interna}

\subsection{Particionamentos fixos e dependência do teste}

As três condições de cada cenário utilizam o mesmo conjunto de teste real fixo, o que melhora a comparabilidade interna. Entretanto, os resultados permanecem condicionados ao particionamento adotado. Um único particionamento pode privilegiar determinadas regiões da distribuição, períodos ou perfis de pacientes. O Air Quality utiliza separação cronológica, adequada para evitar a mistura ingênua entre passado e futuro, mas ainda representa apenas um corte temporal. Occupancy e SUPPORT2 também podem apresentar variação sob outros particionamentos.

Uma extensão desejável é repetir o protocolo com múltiplos particionamentos válidos para a estrutura do dataset. Essa repetição não deve recorrer automaticamente à validação aleatória em dados temporais ou agrupados, pois tal procedimento pode causar vazamento de informação e produzir estimativas otimistas \cite{Kaufman2012,Roberts2017,Bergmeir2012}. O objetivo seria medir a sensibilidade da preservação ao particionamento sem violar a unidade amostral de cada domínio.

\subsection{Incerteza Monte Carlo}

O critério de parada controla a meia-largura máxima do intervalo de confiança das perdas médias estimadas, mas $\rankrho$, $\AW$, $\AR$ e $\SR$ são funções derivadas dessas perdas, não alvos diretos do critério de parada. Pequenas alterações nas perdas podem mudar rankings, o vencedor efetivo $\Wmat$ e as inversões registradas em $\Rmat$ mesmo depois de as médias atingirem a precisão estabelecida. Empates exatos de perda são tratados de forma determinística por propagação do vencedor anterior após a inicialização (Seção~\ref{sec:winner-relation}); essa regra remove ambiguidade na definição de $\Wmat$, mas não elimina a sensibilidade de $\Wmat$ e $\Rmat$ a deslocamentos pequenos na perda estimada próximos de um cruzamento.

Amostras aninhadas e sementes controladas reduzem o ruído das comparações, porém também induzem dependência entre pontos de $n$ e entre condições experimentais. A dependência é intencional e melhora a comparação pareada, mas impede interpretar as curvas como séries de observações independentes. Trabalhos futuros podem estimar intervalos de confiança dos próprios indicadores de preservação por reamostragem das replicações.

\subsection{Balanceamento das amostras de treinamento}

O protocolo usa $n$ amostras por classe para isolar o efeito do tamanho amostral e evitar que a classe majoritária domine o treinamento. Essa escolha não reproduz necessariamente a prevalência natural dos datasets, sobretudo no SUPPORT2. A literatura de dados desbalanceados mostra que balancear pode melhorar determinados objetivos, mas também altera probabilidades posteriores, fronteiras de decisão e interpretação operacional \cite{He2009}.

O teste permanece real e fixo, mas o treinamento balanceado representa um experimento controlado, não a rotina de implantação mais provável. Portanto, as perdas não devem ser interpretadas diretamente como desempenho esperado em uma população clínica ou operacional sem recalibração de prevalência e custos de erro.

\subsection{Parametrização dos geradores}

A Gaussiana única e o GMM são ajustados a partir dos dados reais de treinamento. A incerteza dos parâmetros, as escolhas de regularização e a seleção do número de componentes influenciam as amostras produzidas e, por consequência, os perfis que emergem após o treinamento dos classificadores nos braços sintéticos. O BIC controla complexidade, mas não garante que o GMM selecionado produza amostras capazes de reproduzir melhor o perfil de cooperação preditiva de referência. A maior fidelidade visual do GMM em alguns cenários também não assegura melhor preservação do perfil.

Os geradores são condicionais à classe e não modelam explicitamente deriva da distribuição, dependência temporal, censura, ausência de dados ou mecanismos de coleta. O protocolo trabalha com uma representação tabular padronizada. A simplificação é adequada ao escopo atual, mas limita a interpretação em processos nos quais a ordem ou a dependência entre observações faz parte do fenômeno.

\subsection{Calibração e aleatoriedade do Random Forest}

O Random Forest foi calibrado uma vez nos dados reais de treinamento do SUPPORT2, e a configuração foi mantida nas três condições. Essa decisão evita favorecer um gerador específico, mas os resultados dependem da grade de calibração e do critério adotado. Outra configuração poderia produzir perdas e estruturas diferentes.

As sementes do estimador são pareadas entre condições por replicação. Esse desenho controla a variabilidade interna do algoritmo e torna a comparação mais precisa, mas pode elevar a concordância em relação a uma execução com sementes independentes. Um estudo de sensibilidade deve comparar políticas pareadas e independentes para separar a preservação atribuível aos dados da redução de variância causada pelo desenho experimental.

\section{Validade de construto}

\subsection{O que os quatro indicadores representam}

A correlação de Spearman mede associação entre rankings completos, não igualdade de perdas. Duas condições podem apresentar $\rankrho$ alto e diferenças relevantes no desempenho absoluto. A \WinnersAgreement{} ($\AW$) mede a proporção de pares com a mesma relação efetiva de vencedor, mas ignora a magnitude da vantagem; pares em que a relação permanece indefinida em qualquer condição são ignorados, não tratados como discordância.

A \ReversalExistenceAgreement{} ($\AR$) é um indicador \textbf{baseado em conjuntos}: ela ignora a direção de cada inversão (recuperável apenas em $\Wmat$) e sua magnitude em perda, respondendo somente se o par inverteu ou não. A \ReversalSampleSizeSimilarity{} ($\SR$) é definida \textbf{apenas no suporte compartilhado de inversões} --- pares que invertem em ambas as condições --- e sua normalização depende explicitamente do prefixo avaliado, $\log_2 n_t-\log_2 n_1$; o mesmo par de tamanhos amostrais absolutos produz valores de $\SR$ diferentes conforme o prefixo em que é avaliado. Além disso, a matriz $\Rmat$ armazena apenas a \textbf{última} inversão de cada par: ela descarta o número total de inversões, sua ordem, sua duração e a magnitude da perda envolvida em cada uma. Dois braços podem compartilhar a última inversão de um par e divergir substancialmente em inversões intermediárias, ou podem discordar sobre a última inversão apesar de exibirem trajetórias de perda praticamente equivalentes.

Os quatro indicadores são, portanto, componentes parciais de um construto mais amplo --- o perfil de cooperação preditiva --- e nenhum deles, isoladamente, o resume. A decisão de não combiná-los em um índice único evita esconder essas diferenças, mas exige interpretação conjunta.

\subsection{Dependência da grade amostral avaliada}

A matriz de inversões $\Rmat(n_t)$ e os indicadores derivados $\AR(n_t)$ e $\SR(n_t)$ dependem explicitamente da grade de tamanhos amostrais avaliada e do prefixo $n_t$ considerado. A última inversão registrada é a mais tardia \emph{dentro da grade}, não necessariamente a última que ocorreria se a grade se estendesse além do maior $n$ avaliado em cada dataset. Um par pode apresentar uma nova inversão além desse limite, invalidando a interpretação de estabilização final.

Também existe uma assimetria entre ausência de inversão e inversão não observada. Uma célula indefinida de $\Rmat$ pode significar que um subconjunto foi superior ao outro durante toda a grade avaliada ou que uma inversão ocorreria além dela. O relatório distingue estados de disponibilidade explicitamente (Seção~\ref{sec:reversal-matrix}), mas não elimina essa ambiguidade substantiva sobre o que existe fora do intervalo observado.

\subsection{Perda 0--1 como medida de desempenho}

A métrica básica é a perda de classificação 0--1. Ela é simples e comparável, porém não incorpora calibração probabilística, gravidade dos erros, sensibilidade, especificidade ou utilidade clínica. No SUPPORT2, por exemplo, dois classificadores com a mesma taxa de erro podem ter perfis de falso positivo e falso negativo muito diferentes.

O perfil de cooperação preditiva pode depender da métrica de desempenho utilizada como base. Trabalhos futuros devem verificar se $\rankrho$, $\AW$, $\AR$ e $\SR$ se mantêm com acurácia balanceada, macro-F1, AUC, perda logarítmica ou funções de custo específicas.

\section{Validade externa}

\subsection{Número e diversidade dos datasets}

O estudo inclui três datasets públicos de domínios distintos, o que melhora a diversidade em relação ao SLACGS idealizado. Ainda assim, três cenários são insuficientes para generalizar relações entre geometria e preservação. Occupancy, Air Quality e SUPPORT2 possuem baixa dimensionalidade, alvo binário e número de atributos compatível com enumeração exaustiva.

Não foram incluídos problemas de alta dimensão, múltiplas classes, imagens, texto, sinais brutos ou séries temporais modeladas explicitamente. Os candidatos CSI Wi-Fi e sistema hidráulico foram adiados justamente porque exigem tratamento de sessões, ciclos e dependência temporal. Assim, os resultados valem para o desenho tabular adotado, não para aprendizado multicanal em geral.

\subsection{Classificadores avaliados}

SVM linear, regressão logística e Gaussian Naive Bayes foram avaliados nos cenários principais. O Random Forest foi executado apenas no SUPPORT2. Portanto, não é possível afirmar que esse algoritmo preserva melhor o perfil de cooperação preditiva em outros datasets. Seu resultado deve ser atribuído à interação entre SUPPORT2, configuração calibrada, gerador e política de sementes.

Modelos não lineares adicionais, como SVM com kernel, boosting, k-NN e redes neurais, podem responder de forma diferente às distribuições sintéticas. A preservação não é propriedade exclusiva do gerador e deve ser reavaliada para cada família de decisão.

\subsection{Interpretação de custos}

O trabalho não coleta dados financeiros nem implementa restrições operacionais reais. Os exemplos de sensores, exames e infraestrutura ilustram a agenda de aplicação, mas não demonstram economia. Custos de aquisição, manutenção, erro e referência variam entre instituições, países e períodos. Uma configuração eficiente em um contexto pode ser inviável em outro.

\section{Limitações computacionais}

A enumeração exaustiva cresce como $2^d-1$. Com cinco atributos, cada cenário possui 31 subconjuntos; com sete, o SUPPORT2 possui 127. O custo total ainda é multiplicado por braços, classificadores, tamanhos amostrais e replicações. Esse crescimento limita a aplicação direta do protocolo a dimensionalidades maiores.

O Random Forest demandou aproximadamente três a cinco vezes o tempo das simulações com classificadores lineares no ambiente utilizado. Por isso, não foi executado em Occupancy e Air Quality dentro do prazo do trabalho. Essa assimetria experimental impede uma comparação cruzada completa, embora o estudo adicional no SUPPORT2 permaneça informativo.

A paralelização por processos reduz o tempo de execução, mas não elimina o custo de milhares de treinamentos de classificadores. Em dimensões maiores, será necessário substituir a enumeração exaustiva por estratégias de busca, amostragem de subconjuntos, ramificação e poda (\textit{branch-and-bound}), métodos de seleção ou análise por cardinalidade. Qualquer aproximação desse tipo deve preservar a capacidade de estudar relações par a par relevantes.

\section{Seleção e apresentação de resultados}

Os relatórios HTML contêm grande volume de gráficos e matrizes. O PDF apresenta uma seleção orientada pelas hipóteses científicas. Essa seleção cria risco de enfatizar casos mais expressivos. Para reduzir esse problema, as tabelas finais incluem todas as combinações auditadas, e os identificadores dos relatórios completos são fornecidos no Apêndice~A.

As visualizações geométricas são exploratórias. A percepção de maior ou menor semelhança com regiões gaussianas é útil para formular hipóteses, mas não substitui testes formais de ajuste ou medidas quantitativas de distância entre distribuições. O relatório evita inferir causalidade a partir dessas figuras.

\section{Limites da extrapolação}

Nenhum resultado demonstra que o perfil preservado até o maior $n$ avaliado continuará válido para amostras maiores. A última inversão registrada em $\Rmat$ é relativa à grade experimental. Projetar $\Rmat$ para $n+a$ exige validação fora da faixa usada para ajustar e selecionar o gerador.

Os resultados tampouco se transferem automaticamente entre domínios. Nem a elevada preservação dos quatro indicadores do Gaussian Naive Bayes no Air Quality, nem o padrão de $\AR$/$\SR$ observado no Random Forest no SUPPORT2 (Seção~\ref{sec:support2-rf}), autorizam generalizações para outros conjuntos de sensores ou classificadores. As conclusões constituem evidência de viabilidade em alguns cenários, não uma validação geral: a intensidade e a utilidade da preservação devem ser novamente demonstradas em cada aplicação.

\section{Escopo e completude dos dados apresentados}

Os quatro indicadores de preservação --- $\rankrho$, $\AW$, $\AR$ e $\SR$, complementados por $\DR$ --- são reportados no Capítulo~5 para os quatro cenários de escala completa (Occupancy 000002, Air Quality 000005, SUPPORT2 000007 e 000008), calculados diretamente pelo código atual (\texttt{coinfosim.results.predictive\_profile}) a partir dos resultados brutos por replicação persistidos para essas execuções. Não há, portanto, indicador reportado como indisponível por ausência de dado bruto neste relatório; as limitações que permanecem em aberto --- particionamento fixo, incerteza Monte Carlo, balanceamento de classes, número de datasets e classificadores, e os limites de extrapolação além da grade avaliada --- são as discutidas nas seções anteriores deste capítulo, e não decorrem de indisponibilidade de dados.
